\chapter{Forward-Rate and Market Models}\label{ch:rc:forward-rate-and-market-models}

For twenty years before 1997, traders priced caps with Black's formula applied
to each forward rate as if it were a stock, knowing that the formula came from
no consistent model of the curve: the lognormal forward of one caplet and the
lognormal forward of the next could not both be lognormal under the same
measure, and \omterm{def:rc:short-rate-models:short}{short-rate models} priced caps with other formulas. In 1997 three
papers, by Brace, Gatarek and Musiela, by Jamshidian, and by Miltersen,
Sandmann and Sondermann, showed that the traders' formula was exactly right in
a model in which every forward rate is lognormal under the measure of its own
payment date, and that one such model could hold all forwards at once. The
market model, as it came to be called, calibrates to caplets by construction,
needs Monte Carlo for almost everything else, and makes visible a parameter the
\omterm{def:rc:short-rate-models:short}{short-rate models} hid: the correlation between forward rates. This chapter
derives the general framework it belongs to, builds and simulates it, and shows
two correlations that fit the same caplets and price the same swaption 22\%
apart.

\section{The Heath--Jarrow--Morton framework and its drift}

\begin{definition}[Heath--Jarrow--Morton framework]\label{def:rc:forward-rate-and-market-models:hjm}
The \emph{Heath--Jarrow--Morton framework}\index{Heath--Jarrow--Morton framework}
(1992) models the whole instantaneous forward curve: for each maturity $T$,
$df(t,T) = \alpha(t,T)\,dt+\sigma_f(t,T)\,dW_t$ under $\mathbb Q$, starting from
today's curve $f(0,T)$, with volatilities $\sigma_f(t,T)$ chosen freely (vectors
for several factors).
\end{definition}

\begin{definition}[HJM drift condition]\label{def:rc:forward-rate-and-market-models:drift}
The \emph{HJM drift condition}\index{HJM drift condition} is the restriction
that absence of arbitrage imposes on the drift once the volatilities are
chosen:
\[
\alpha(t,T) = \sigma_f(t,T)\int_t^T\sigma_f(t,u)\,du .
\]
\end{definition}

\begin{proposition}[Why the drift is fixed]\label{prop:rc:forward-rate-and-market-models:drift}
Bond prices $P(t,T) = \exp(-\int_t^Tf(t,u)\,du)$ discounted by the bank account
are $\mathbb Q$-martingales if and only if the drift condition holds.
\end{proposition}
\begin{proof}
$d\ln P(t,T) = f(t,t)\,dt-\int_t^Td f(t,u)\,du = \bigl(r_t-\int_t^T\alpha(t,u)\,du\bigr)dt
-\Sigma(t,T)\,dW_t$ with $\Sigma(t,T)=\int_t^T\sigma_f(t,u)\,du$. By Itô's formula
the drift of $P/B$ is $-\int_t^T\alpha\,du+\frac12\Sigma^2$; it vanishes for all
$T$ iff $\int_t^T\alpha(t,u)\,du = \frac12\Sigma(t,T)^2$, and differentiating in $T$
gives the condition.
\end{proof}

\begin{example}[Hull--White is an HJM model]\label{ex:rc:forward-rate-and-market-models:hw}
Chapter~7's model has $\sigma_f(t,T) = \sigma e^{-\kappa(T-t)}$: its drift is
$\sigma^2e^{-\kappa(T-t)}(1-e^{-\kappa(T-t)})/\kappa$, 4.1 basis points a year for the
ten-year forward with $\sigma = 80$ basis points and $\kappa = 3\%$. Volatilities of
this exponential form are exactly those for which the forward curve is driven by
a single Markov state; a general $\sigma_f$ makes the short rate path-dependent
and the model can only be simulated.
\end{example}

\section{The market model}

The instantaneous forward is not traded. The market model takes the forwards
that are: $F_k(t)$ for the accrual periods $[T_k,T_{k+1}]$ of a tenor structure.

\begin{definition}[LIBOR market model]\label{def:rc:forward-rate-and-market-models:lmm}
The \emph{LIBOR market model}\index{LIBOR market model} (or lognormal forward
market model) makes each forward rate lognormal with a deterministic volatility
under the forward measure of its payment date:
$dF_k(t) = \sigma_k(t)F_k(t)\,dW^{T_{k+1}}_k(t)$, with $d\langle W_i,W_j\rangle =
\rho_{ij}\,dt$. Its name comes from the interbank index it was built for; the
dynamics serve any tenor structure of forward-looking rates, and chapter~10
extends them to overnight-rate compounding.
\end{definition}

Each caplet is then priced exactly by Black's formula with volatility
$\sqrt{\frac1{T_k}\int_0^{T_k}\sigma_k(t)^2dt}$, which is what makes the model
calibrate to caps by construction. A common volatility form is Rebonato's
\[
\sigma_k(t) = \phi_k\,g(T_k-t),\qquad g(\tau) = (a+b\tau)e^{-c\tau}+d ,
\]
a hump in the time to reset scaled by a constant per forward
(\cref{fig:rc:forward-rate-and-market-models:abcd}); the $\phi_k$ are solved so
that every caplet reprices.

\begin{omfigure}
\begin{tikzpicture}
\begin{axis}[width=11cm,height=5cm,
  xlabel={time to reset $\tau$ (years)},ylabel={$g(\tau)$},
  tick label style={font=\small},label style={font=\small},
  xmin=0,xmax=10,ymin=0,ymax=0.35,grid=major,grid style={gray!25},
  table/col sep=comma]
\addplot[omThm,very thick] table[x=tau,y=g]{figdata/rates-credit-risk/08-forward-rate-and-market-models/abcd.csv};
\end{axis}
\end{tikzpicture}

\omcaption{Rebonato's volatility function with $a=0.05$, $b=0.10$, $c=0.60$,
$d=0.15$: a forward's volatility is highest about a year and a half before its
reset. Scaled per forward, it makes the volatility of each forward depend on
its time to reset, so the volatility curve seen in a year looks like today's.
Data: the chapter's tutorial.}\label{fig:rc:forward-rate-and-market-models:abcd}
\end{omfigure}

\section{Drifts under the spot and terminal measures}

To simulate all forwards together they must share one measure, and under any
single measure all but one of them acquire a drift.

\begin{definition}[Spot measure, terminal measure]\label{def:rc:forward-rate-and-market-models:measures}
The \emph{spot measure}\index{spot measure} $\mathbb Q^d$ of a tenor structure has
as numeraire the discretely rebalanced bank account $B_d(t) =
P(t,T_{m(t)})\prod_{j<m(t)}(1+\delta_jF_j(T_j))$, which holds the bond maturing at
the next reset date $T_{m(t)}$ and rolls into the next one at each reset. The
\emph{terminal measure}\index{terminal measure} is the forward measure of the last
payment date $T_n$, with numeraire $P(t,T_n)$.
\end{definition}

\begin{proposition}[Market-model drifts]\label{prop:rc:forward-rate-and-market-models:spotdrift}
Under the \omterm{def:rc:forward-rate-and-market-models:measures}{spot measure}, for a forward $F_k$ alive at $t$,
\[
\frac{dF_k}{F_k} = \sigma_k(t)\sum_{j=m(t)}^{k}\frac{\delta_j\rho_{kj}\sigma_j(t)F_j}{1+\delta_jF_j}\,dt+\sigma_k(t)\,dW_k,
\]
and under the \omterm{def:rc:forward-rate-and-market-models:measures}{terminal measure} the drift is $-\sigma_k\sum_{j=k+1}^{n-1}
\delta_j\rho_{kj}\sigma_jF_j/(1+\delta_jF_j)$.
\end{proposition}
\begin{proof}[Partial proof]
Change numeraire one period at a time: the density of $\mathbb Q^{T_{j+1}}$ with
respect to $\mathbb Q^{T_j}$ is proportional to $P(t,T_{j+1})/P(t,T_j) =
1/(1+\delta_jF_j)$, whose volatility is $-\delta_j\sigma_jF_j/(1+\delta_jF_j)$; by
Girsanov's theorem (One Quant Book 4, chapter 5) each change adds the
covariance of $\ln F_k$ with that density to the drift. Summing from the
numeraire's bond to $F_k$'s own payment date gives both formulas.
\end{proof}

\begin{omfigure}
\begin{tikzpicture}[x=1.25cm,font=\small]
\draw[->,thick] (0,0) -- (9.6,0) node[right] {time};
\foreach \x in {0,...,8} {\draw (\x,-0.1) -- (\x,0.1); \node[below,font=\footnotesize] at (\x,-0.12) {$T_{\x}$};}
\foreach \x in {0,...,7} {\draw[omDef,fill=omDef!10] (\x+0.06,0.35) rectangle (\x+0.94,0.75); \node[font=\scriptsize] at (\x+0.5,0.55) {$F_{\x}$};}
\draw[omThm,very thick,-{Stealth[length=2mm]}] (0,1.25) -- (1,1.25) -- (1,1.55) -- (2,1.55) -- (2,1.85) -- (3,1.85);
\node[text=omThm,anchor=west,font=\footnotesize] at (3.1,1.85) {spot measure: roll the next bond at each reset};
\draw[omProp,very thick,-{Stealth[length=2mm]}] (0,2.6) -- (8,2.6);
\node[text=omProp,anchor=south,font=\footnotesize] at (4,2.65) {terminal measure: the bond maturing at the last date};
\end{tikzpicture}

\omcaption{A tenor structure of eight annual forwards and the numeraires of the
two common simulation measures. Under the \omterm{def:rc:forward-rate-and-market-models:measures}{spot measure} each forward's drift
sums over the forwards between the next reset and itself; under the \omterm{def:rc:forward-rate-and-market-models:measures}{terminal
measure}, over those after it.}\label{fig:rc:forward-rate-and-market-models:tenor}
\end{omfigure}

Simulation steps $\ln F_k$ with the drift evaluated at the start of the step
(Euler) and again at the predicted end, averaging the two
(predictor--corrector): the drift depends on the forwards themselves, and the
correction keeps the simulated caplets on their Black prices at coarse steps.

\omcode{code/firm/lmm/firm_lmm.py}{86}{114}{The spot-measure drift and the predictor--corrector step.}\label{lst:rc:forward-rate-and-market-models:step}

\begin{example}[Caplets come back]\label{ex:rc:forward-rate-and-market-models:caplets}
Ten annual forwards from chapter~2's euro OIS curve (1.98\% for the first year
rising to 3.03\% for the tenth), caplet volatilities of 70 basis points normal
converted to Black (33.0\% at one year down to 23.1\% at nine), Rebonato's
function with the parameters of \cref{fig:rc:forward-rate-and-market-models:abcd}
and correlation $\rho_{ij}=e^{-0.1|i-j|}$: 20\,000 paths under the \omterm{def:rc:forward-rate-and-market-models:measures}{spot measure}, four
steps a year, reprice the at-the-money caplets to within two standard errors
(\cref{fig:rc:forward-rate-and-market-models:caplets}).
\end{example}

\begin{omfigure}
\begin{tikzpicture}
\begin{axis}[width=11cm,height=5.2cm,
  xlabel={caplet expiry (years)},ylabel={Black volatility (\%)},
  tick label style={font=\small},label style={font=\small},
  xmin=0.5,xmax=9.5,ymin=20,ymax=36,grid=major,grid style={gray!25},
  legend columns=2,legend style={font=\footnotesize,at={(0.5,-0.25)},anchor=north,draw=none,fill=none,/tikz/every even column/.append style={column sep=8pt}},
  table/col sep=comma]
\addplot[omDef,thick] table[x=k,y=target]{figdata/rates-credit-risk/08-forward-rate-and-market-models/caplets.csv};
\addplot[omThm,only marks,mark=*,mark size=2pt,error bars/.cd,y dir=both,y explicit] table[x=k,y=mc,y error=err]{figdata/rates-credit-risk/08-forward-rate-and-market-models/caplets.csv};
\legend{target caplet volatility,Monte Carlo ($\pm2$ standard errors)}
\end{axis}
\end{tikzpicture}

\omcaption{At-the-money caplet volatilities implied from Monte Carlo prices of
the calibrated market model, against the targets. Data: the chapter's
tutorial.}\label{fig:rc:forward-rate-and-market-models:caplets}
\end{omfigure}

\section{Calibrating to caps and swaptions: Rebonato's formula}

A swap rate is a weighted average of forwards,
$S_{a,b}=\sum_{i=a}^{b-1}w_iF_i$ with $w_i = \delta_iP(t,T_{i+1})/A_{a,b}(t)$, so its
volatility follows from theirs and their correlations.

\begin{definition}[Swap market model]\label{def:rc:forward-rate-and-market-models:smm}
A \emph{swap market model}\index{swap market model} makes chosen swap rates
lognormal under their own annuity measures, so that swaptions are priced by
Black's formula. It cannot hold together with a lognormal market model of the
forwards: a weighted sum of lognormals is not lognormal. Desks use one and
approximate the other.
\end{definition}

\begin{definition}[Rebonato's formula]\label{def:rc:forward-rate-and-market-models:rebonato}
\emph{Rebonato's formula}\index{Rebonato's formula} approximates the Black
volatility of the swaption expiring at $T_a$ on $S_{a,b}$ in the market model by
freezing the weights and the forwards at today's values:
\[
\sigma_{a,b}^2\,T_a \approx \frac1{S_{a,b}(0)^2}\sum_{i,j=a}^{b-1}
w_i(0)w_j(0)F_i(0)F_j(0)\,\rho_{ij}\int_0^{T_a}\sigma_i(t)\sigma_j(t)\,dt .
\]
\end{definition}

\omcode{code/firm/lmm/firm_lmm.py}{75}{83}{Rebonato's frozen-weights swaption volatility.}\label{lst:rc:forward-rate-and-market-models:rebonato}

\section{Correlation and factor reduction}

\begin{definition}[Forward-rate correlation]\label{def:rc:forward-rate-and-market-models:corr}
The \emph{forward-rate correlation}\index{forward-rate correlation} $\rho_{ij}$ of
a market model is the instantaneous correlation of the Brownian motions
driving $F_i$ and $F_j$. It is usually parametrised, for example
$\rho_{ij} = e^{-\beta|T_i-T_j|}$, and reduced to two or three factors by keeping
the leading eigenvectors of the matrix (the principal components of
chapter~3), which speeds simulation and removes noise.
\end{definition}

Caplets do not depend on correlation; swaptions do, since a swap rate averages
several forwards and averaging decorrelated rates reduces volatility. A market
model calibrated to caplets alone therefore leaves its swaption prices to an
assumption.

\begin{example}[Same caplets, different swaptions]\label{ex:rc:forward-rate-and-market-models:corr}
Calibrate the model to the same caplets with $\beta=0.02$ (forwards nearly
perfectly correlated) and with $\beta=0.40$ (the one-year and nine-year forwards
correlated at 0.04). The swaption expiring in five years into one year is
priced at 25.0\% of Black volatility by both (it is one forward); into five years,
at 22.1\% and 17.3\% (\cref{fig:rc:forward-rate-and-market-models:swaptions}). On
EUR~100 million, at the money (forward 2.933\%, annuity 4.102), that is EUR~2.35
million against 1.84 million: 22\% apart for the same caplets. Monte Carlo
confirms \omterm{def:rc:forward-rate-and-market-models:rebonato}{Rebonato's formula} to within its standard errors.
\end{example}

\begin{omfigure}
\begin{tikzpicture}
\begin{axis}[width=11cm,height=5cm,
  xlabel={forward index $j$},ylabel={$\rho_{1j}$},
  tick label style={font=\small},label style={font=\small},
  xmin=1,xmax=9,ymin=0,ymax=1.05,grid=major,grid style={gray!25},
  legend columns=3,legend style={font=\footnotesize,at={(0.5,-0.25)},anchor=north,draw=none,fill=none,/tikz/every even column/.append style={column sep=8pt}},
  table/col sep=comma]
\addplot[omDef,thick,mark=*] table[x=j,y=b02]{figdata/rates-credit-risk/08-forward-rate-and-market-models/corr.csv};
\addplot[omProp,thick,mark=triangle*] table[x=j,y=b10]{figdata/rates-credit-risk/08-forward-rate-and-market-models/corr.csv};
\addplot[omThm,thick,mark=square*] table[x=j,y=b40]{figdata/rates-credit-risk/08-forward-rate-and-market-models/corr.csv};
\legend{$\beta=0.02$,$\beta=0.10$,$\beta=0.40$}
\end{axis}
\end{tikzpicture}

\omcaption{Correlation of the first annual forward with the others under the
exponential parametrisation $\rho_{ij}=e^{-\beta|T_i-T_j|}$ for three values of
$\beta$. Data: the chapter's tutorial.}\label{fig:rc:forward-rate-and-market-models:corrplot}
\end{omfigure}

\begin{omfigure}
\begin{tikzpicture}
\begin{axis}[width=11.5cm,height=5.4cm,
  xlabel={tenor of the swap starting in five years (years)},ylabel={Black volatility (\%)},
  tick label style={font=\small},label style={font=\small},
  xmin=0.8,xmax=5.2,ymin=15,ymax=27,xtick={1,2,3,4,5},grid=major,grid style={gray!25},
  legend columns=2,legend style={font=\footnotesize,at={(0.5,-0.25)},anchor=north,draw=none,fill=none,/tikz/every even column/.append style={column sep=8pt}},
  table/col sep=comma]
\addplot[omDef,thick] table[x=tenor,y=reb02]{figdata/rates-credit-risk/08-forward-rate-and-market-models/swaptions.csv};
\addplot[omDef,only marks,mark=*] table[x=tenor,y=mc02]{figdata/rates-credit-risk/08-forward-rate-and-market-models/swaptions.csv};
\addplot[omThm,thick,dashed] table[x=tenor,y=reb40]{figdata/rates-credit-risk/08-forward-rate-and-market-models/swaptions.csv};
\addplot[omThm,only marks,mark=square*] table[x=tenor,y=mc40]{figdata/rates-credit-risk/08-forward-rate-and-market-models/swaptions.csv};
\legend{Rebonato $\beta=0.02$,Monte Carlo $\beta=0.02$,Rebonato $\beta=0.40$,Monte Carlo $\beta=0.40$}
\end{axis}
\end{tikzpicture}

\omcaption{Five-year-expiry swaption volatilities from two market models that
reprice the same caplets. Decorrelating the forwards lowers the volatility of
longer swap rates; \omterm{def:rc:forward-rate-and-market-models:rebonato}{Rebonato's formula} (lines) tracks Monte Carlo (marks).
Data: the chapter's tutorial.}\label{fig:rc:forward-rate-and-market-models:swaptions}
\end{omfigure}

\begin{remark}[Calibrating the correlation]\label{rem:rc:forward-rate-and-market-models:calib}
Desks calibrate the market model jointly: the $\phi_k$ to caplets or co-terminal
swaptions, the volatility shape and the correlation to the \omterm{def:rc:sabr-in-rates-and-the-volatility-cube:matrix}{swaption matrix},
often by fitting $\beta$ (or a two-parameter form) to historical \omterm{def:rc:forward-rate-and-market-models:corr}{forward-rate
correlations} and letting the swaptions fix the rest. A product depending on the
correlation of rates that no liquid instrument pins down carries a reserve
for it (\cref{ch:rc:pnl-explain-and-independent-price-verification}).
\end{remark}

\section{Tutorial: a market model on the euro curve}\label{tut:rc:forward-rate-and-market-models:tutorial}

\begin{tutorial}
\textbf{Goal.} Calibrate a market model to caplets, simulate it, and measure how
correlation moves swaptions. \textbf{End state:}
\cref{fig:rc:forward-rate-and-market-models:caplets,fig:rc:forward-rate-and-market-models:swaptions}.
\begin{tutsteps}
\item \textbf{Forwards and targets}: \texttt{F0} from chapter~2's €STR curve, caplet
Black volatilities from 70 basis points normal.
\item \textbf{Calibrate}: \texttt{calibrate\_to\_caplets(F0, vols, abcd, beta)}
solves the $\phi_k$.
\item \textbf{Simulate}: \texttt{mc\_caplet} and \texttt{mc\_swaption} under the \omterm{def:rc:forward-rate-and-market-models:measures}{spot
measure}; \texttt{caplet\_check()}.
\item \textbf{Correlation}: \texttt{swaption\_table()} for $\beta=0.02$ and $0.40$;
\texttt{fig\_rc\_lmm.py} writes the charts.
\end{tutsteps}
\textbf{What to change next.} Simulate under the \omterm{def:rc:forward-rate-and-market-models:measures}{terminal measure} and compare
standard errors; reduce the correlation matrix to two factors and measure the
change in the five-into-five volatility.
\end{tutorial}

\section{Build: a market-model engine}\label{bld:rc:forward-rate-and-market-models:lmm}

\begin{build}
\bfield{Purpose} The benchmark model for products that depend on several rates
at several dates: Bermudans (chapter~9), \omterm{def:rc:convexity-adjustments-and-constant-maturity-products:spread}{CMS spread options} (chapter~6), and the
validation of chapter~7's one-factor model.
\bfield{Interface} \texttt{LMM(f0, phi, abcd, beta)} with \texttt{vol},
\texttt{corr}, \texttt{integrated\_cov}, \texttt{discount\_factors}, \texttt{swap};
\texttt{calibrate\_to\_caplets}; \texttt{rebonato\_vol}; \texttt{simulate},
\texttt{mc\_caplet}, \texttt{mc\_swaption}; \texttt{hjm\_drift\_gaussian}.
\bfield{Rules} Annual tenor structure, lognormal forwards, \omterm{def:rc:forward-rate-and-market-models:measures}{spot measure} with
predictor--corrector steps; fixed seeds; numpy only.
\bfield{Acceptance tests} \texttt{code/firm/lmm/tests/}: the calibrated integrated
variances match the caplet targets and \omterm{def:rc:forward-rate-and-market-models:rebonato}{Rebonato's formula} for a one-period
swaption returns the \omterm{def:rc:vanilla-rates-options:flat}{caplet volatility}; Monte Carlo reprices a caplet within
three standard errors; a deep in-the-money swaption is the forward swap; lower
correlation lowers swaption volatility; the Hull--White HJM drift has its
closed form.
\bfield{Stretch} Shifted or normal forwards for low rates; a stochastic-volatility
market model; calibration of $\beta$ to the \omterm{def:rc:sabr-in-rates-and-the-volatility-cube:matrix}{swaption matrix}.
\end{build}

\begin{omsources}
\item D.\ Heath, R.\ Jarrow and A.\ Morton, ``Bond pricing and the term structure of
interest rates'', \emph{Econometrica} 60(1), 1992.
\item A.\ Brace, D.\ Gatarek and M.\ Musiela, ``The market model of interest rate
dynamics'', \emph{Mathematical Finance} 7, 1997; F.\ Jamshidian, ``LIBOR and swap
market models and measures'', \emph{Finance and Stochastics} 1, 1997; K.\
Miltersen, K.\ Sandmann and D.\ Sondermann, \emph{Journal of Finance} 52(1), 1997.
\item R.\ Rebonato, \emph{Modern Pricing of Interest-Rate Derivatives: The LIBOR
Market Model and Beyond}, Princeton University Press, 2002.
\end{omsources}

\section{Exercises}

\begin{exercise}[$\star$]\label{exo:rc:forward-rate-and-market-models:1}
In an HJM model with $\sigma_f(t,T)=\sigma$ constant, what is the drift of $f(t,T)$?
Which \omterm{def:rc:short-rate-models:short}{short-rate model} is it?
\end{exercise}

\begin{exercise}[$\star$]\label{exo:rc:forward-rate-and-market-models:2}
Why is each caplet priced exactly by Black's formula in the market model, but
not each swaption?
\end{exercise}

\begin{exercise}[$\star$]\label{exo:rc:forward-rate-and-market-models:3}
With $\beta = 0.40$, what is the correlation between the first and the ninth annual
forwards? With $\beta=0.02$?
\end{exercise}

\begin{exercise}[$\star\star$]\label{exo:rc:forward-rate-and-market-models:4}
Under the \omterm{def:rc:forward-rate-and-market-models:measures}{spot measure}, write the drift of $F_2$ at $t=0.5$ in a model with annual
forwards. Which forwards appear?
\end{exercise}

\begin{exercise}[$\star\star$]\label{exo:rc:forward-rate-and-market-models:5}
Two forwards of 3\% with Black volatility 20\% each, weights 0.5 each and
correlation $\rho$. With \omterm{def:rc:forward-rate-and-market-models:rebonato}{Rebonato's formula}, what is the swap-rate volatility for
$\rho=1$ and $\rho=0.5$?
\end{exercise}

\begin{exercise}[$\star\star$]\label{exo:rc:forward-rate-and-market-models:6}
Read \cref{fig:rc:forward-rate-and-market-models:swaptions}: why do both models
agree on the five-into-one swaption?
\end{exercise}

\begin{exercise}[$\star\star\star$]\label{exo:rc:forward-rate-and-market-models:7}
\emph{Coding.} Price the five-into-five at-the-money payer on EUR~100 million under
both correlations with \omterm{def:rc:forward-rate-and-market-models:rebonato}{Rebonato's formula}, and give the difference.
\end{exercise}

\begin{exercise}[$\star\star\star$]\label{exo:rc:forward-rate-and-market-models:8}
\emph{Find the flaw.} ``We simulate the market model under the \omterm{def:rc:forward-rate-and-market-models:measures}{spot measure} with
each forward's drift set to zero, because each forward is a martingale.''
\end{exercise}

\section{Problem: The Correlation You Cannot See}

\begin{problem}[{Weekend problem --- a market model calibrated to caplets}]\label{pb:rc:forward-rate-and-market-models:1}
A desk runs a ten-forward annual market model on the euro OIS curve, calibrated
every morning to caplets at 70 basis points of normal volatility. A client
asks for the five-into-five at-the-money payer swaption on EUR~100 million. The
desk's correlation parameter $\beta$ has never been calibrated.

\medskip\noindent\textbf{Part I --- What the caplets fix.}
\begin{enumerate}
\item Give the forward swap rate and the annuity.
\item Which parameters do the caplets determine, and which not?
\item Give the caplet Black volatilities at one and nine years.
\item Why is the five-into-one swaption the same under any correlation?
\item What else is left free by the calibration?
\end{enumerate}

\medskip\noindent\textbf{Part II --- Two correlations.}
\begin{enumerate}[resume]
\item Give the five-into-five volatility with $\beta = 0.02$ and with $\beta=0.40$.
\item Give the two premiums.
\item Give the relative difference.
\item How well does \omterm{def:rc:forward-rate-and-market-models:rebonato}{Rebonato's formula} agree with Monte Carlo here?
\item Which correlation would you expect from historical data, and why?
\end{enumerate}

\medskip\noindent\textbf{Part III --- Fixing it.}
\begin{enumerate}[resume]
\item Which instruments would pin down the correlation?
\item If the market quotes the five-into-five at 20\%, what $\beta$ does it imply
(bracket it)?
\item After calibrating $\beta$ to it, what would still be unpinned?
\item How does the choice of $\beta$ change the desk's hedges?
\item Why is a one-factor \omterm{def:rc:short-rate-models:hw}{Hull--White model} equivalent to $\beta=0$?
\end{enumerate}

\medskip\noindent\textbf{Part IV --- Judgement.}
\begin{enumerate}[resume]
\item Which price would you quote, and what reserve would you hold?
\item How would you monitor the parameter over time?
\item What should the model documentation say about $\beta$?
\item State the \emph{named result}: the five-into-five premiums under the two
correlations and their difference.
\item In one sentence: what does a swaption know that a caplet does not?
\end{enumerate}
\end{problem}

\section{Interview questions}

\begin{interviewq}[$\star$ \iqroles{researcher}]\label{iq:rc:forward-rate-and-market-models:1}
State the \omterm{def:rc:forward-rate-and-market-models:drift}{HJM drift condition} and explain where it comes from.
\end{interviewq}

\begin{interviewq}[$\star\star$ \iqroles{researcher, developer}]\label{iq:rc:forward-rate-and-market-models:2}
How do you simulate a \omterm{def:rc:forward-rate-and-market-models:lmm}{LIBOR market model}, and why is there a drift?
\end{interviewq}

\begin{interviewq}[$\star\star$ \iqroles{researcher, trader}]\label{iq:rc:forward-rate-and-market-models:3}
Why can caplets and swaptions not both be exactly lognormal?
\end{interviewq}

\begin{interviewq}[$\star\star$ \iqroles{trader}]\label{iq:rc:forward-rate-and-market-models:4}
You calibrate a market model to caplets only. What is wrong with your swaption
prices?
\end{interviewq}

\begin{interviewq}[$\star\star\star$ \iqroles{researcher, risk}]\label{iq:rc:forward-rate-and-market-models:5}
Compare Hull--White and the market model for pricing a ten-year Bermudan.
\end{interviewq}

\begin{interviewq}[$\star\star\star$ \iqroles{developer}]\label{iq:rc:forward-rate-and-market-models:6}
Your simulated caplet prices drift away from Black as you use coarser time
steps. What is happening, and what do you change?
\end{interviewq}
